% ---
% titre: Problème de nombres relatifs - la vie de Platon
% theme: NO
% chapitre: Nombres relatifs
% sous_chapitre: Addition et soustraction
% niveau: [10e]
% type: EVAL
% tags: [c-addition,c-soustraction,c-nombres-relatifs,p-question-TS]
% difficulte: 2
% corrige: true
% ---

\question[4\half]
Platon naquit à Athènes en 427 avant Jésus-Christ. Il avait 28 ans lorsque son maître Socrate mourut. En 377 av. J.-C., Platon fonda une école de philosophie, l'Académie, dans laquelle il enseigna jusqu'à sa mort, survenue en 348 av. J.-C.


\begin{parts}
\vspace{20pt}\part En quelle année est mort Socrate ?

\vspace{10pt}{\footnotesize\raggedleft Espace pour ta démarche et tes calculs\par}
\begin{solutionorgrid}[80pt]
-427 + 28 = -399 \textbf{(0.5pt pour le calcul, 0.5pt pour la réponse)}
\end{solutionorgrid}

\begin{tcolorbox}[enhanced jigsaw, interior hidden, colframe=box, parbox=false]%
Réponse: 
\sol{\vspace{20pt}}{Socrate est mort en 399 av J.-C. \textbf{(0.5pt)}}
\end{tcolorbox}


\vspace{20pt}\part Quel âge avait Platon lorsqu'il fonda l'Académie ?

\vspace{10pt}{\footnotesize\raggedleft Espace pour ta démarche et tes calculs\par}
\begin{solutionorgrid}[80pt]
-377 - (-427) = 50 \textbf{(0.5pt pour le calcul, 0.5pt pour la réponse)}
\end{solutionorgrid}

\begin{tcolorbox}[enhanced jigsaw, interior hidden, colframe=box, parbox=false]%
Réponse: 
\sol{\vspace{20pt}}{Il avait 50 ans lorsqu'il fonda l'Académie. \textbf{(0.5pt)}}
\end{tcolorbox}

\vspace{20pt}\part À quel âge mourut Platon?

\vspace{10pt}{\footnotesize\raggedleft Espace pour ta démarche et tes calculs\par}
\begin{solutionorgrid}[80pt]
-348 - (-427) = 79 \textbf{(0.5pt pour le calcul, 0.5pt pour la réponse)}
\end{solutionorgrid}

\begin{tcolorbox}[enhanced jigsaw, interior hidden, colframe=box, parbox=false]%
Réponse: 
\sol{\vspace{20pt}}{Il mourut à 79 ans. \textbf{(0.5pt)}}
\end{tcolorbox}


\end{parts}



